\section{Appendix: Scope and Qualifications}
This concise study compares the two interpretations developed in
\work{All Saints: Communion, Conversion, and the Promised Kingdom},
the expansive study under the same research owner. It concerns the
Vatican typical edition of the 1962 Roman Missal, General Roman Calendar,
on 1 November 2026, with the three prayers of the privileged Sunday
commemoration. The distinct votive Mass, local calendars, present
permission to use the book, and territorial holy-day discipline are
outside its scope. The occurrence, Creed, and Common Preface follow
MR 1962, pp.~702--703, \work{Rubricae generales} 91, 106,
108--109, 111, 113, and \work{Rubricae generales Missalis romani}
475, 480, 482--483, 494(b), 498, 505.

The complete appointed texts and historical study English are printed
in the expansive study. This companion selects original commentary;
it supplies no substitute prayer translation. Its two interpretations
are editorial syntheses. Each Father is cited for his own passage
or question. Schuster supplies reception of the feast and the three
shared Sunday prayers; his historical rubrics and Preface do not
determine the dated celebration. The historical English sources
and native versus translated Psalm numbering remain distinguished.
No fresh Greek or critical Latin collation, exhaustive reception
claim, or independent Breviary verification is made.

The shared historical appendix retains the full chronology evidence,
including the alternatives behind its compact table. Research scope,
source bindings, interpretation reasoning, and textual inspection are
recorded in \texttt{research/scope.md},
\texttt{research/source-bindings.toml},
\texttt{research/interpretations.md}, and \texttt{propers/verified.md}.
Author proof checks are recorded separately from independent reviews.
Historical source wording retains its recorded United States
public-domain basis; modern source editions and apparatus are not
redistributed whole. Original commentary is covered by the project
license, with third-party matter excluded as the final notice states.

\section{References}
\begin{itemize}
\item \work{Missale Romanum}, Vatican editio typica (1962),
All Saints, pp.~702--703; Sunday orations, pp.~416--417;
Common Preface, pp.~252--253; General Rubrics and General
Rubrics of the Missal at the numbers stated above.
\href{https://media.churchmusicassociation.org/pdf/missale62.pdf}{CMAA facsimile}.
\item \work{Douay--Rheims Bible}, Challoner revision,
\href{https://www.gutenberg.org/ebooks/1581}{Gutenberg 1581}:
Ps.~32:1; 33:10--11; Wis.~2:12--3:9;
Matt.~4:23--5:12; 11:25--30; Apoc.~1:9; 6:12--7:12;
Gal.~3:29. Vulgate Psalm numbering first.
\item Bede, \work{The Explanation of the Apocalypse},
trans. Edward Marshall (James Parker, 1878), I, on
Apoc.~7:4--12, pp.~46--53.
\item Augustine, \work{Our Lord's Sermon on the Mount},
I.1.3--I.4.12, Findlay's historical English with
Schaff's revision and notes, NPNF first series 6,
retained CCEL text. Editorial notes are separate from Augustine.
\item Augustine, \work{Expositions of the Psalms},
Ps.~32, enarratio I, \S1; Ps.~33, enarratio II,
\S\S14--15; Ps.~87, \S5. NPNF first series 8
numbers these Psalm XXXIII \S1, Psalm XXXIV
\S\S13--14, and Psalm LXXXVIII \S4.
\item Augustine, \work{The City of God}, XX.9,
Dods edition, II (1871), Gutenberg historical English.
\item John Chrysostom, \work{Homilies on Matthew},
XV.5--7 and XXXVIII.3--5, NPNF first series 10,
historical English revised by Prevost and Riddle,
retained CCEL text.
\item Thomas Aquinas, \work{Super Psalmos}, Ps.~33,
nn.~10--11, \href{https://www.corpusthomisticum.org/cps31.html}{Corpus Thomisticum};
Latin paraphrased.
\item Gregory the Great, \work{Moralia in Iob},
XII.5.6--7, on Job~14:8--9, reusing Wis.~3:2;
\work{Library of the Fathers} historical English,
registered Lectionary Central transcription.
\item Alfredo Ildefonso Schuster, \work{The Sacramentary},
III (Burns Oates \& Washbourne, 1927), pp.~188--189;
V (1930), pp.~208--211.
\end{itemize}
The chronology's source particulars and qualifications are cited
in the same owner's historical appendix: Corbett, ``King David'';
Gigot, ``Book of Wisdom''; Jacquier, ``Gospel of St.~Matthew'';
Durand, ``The New Testament''; van den Biesen, ``Apocalypse''
(\work{Catholic Encyclopedia}); the USCCB Psalms introduction;
and the Haydock facsimile at the loci specified there.
